\documentclass[twoside,12pt]{book}
\usepackage[utf8]{inputenc}
\usepackage[T1]{fontenc}
\usepackage{mathpazo}
\usepackage{amsmath,amssymb,amsthm}
\usepackage{graphicx}
\usepackage{listings}
\usepackage{xcolor}
\usepackage{hyperref}
\usepackage{enumitem}
\usepackage{tikz}
\usepackage{algorithm}
\usepackage{algpseudocode}
\usepackage{booktabs}
\usepackage{geometry}
\usepackage{fancyhdr}
\usepackage{titlesec}
\geometry{papersize={6in,9in},margin=0.75in,top=0.75in,bottom=0.75in,inner=0.9in,outer=0.65in,headheight=14pt}
\pagestyle{fancy}
\fancyhf{}
\fancyhead[LE]{\small\textit{13. Online Algorithms}}
\fancyhead[RO]{\small\textit{Randomized Algorithms}}
\fancyfoot[C]{\thepage}
\renewcommand{\headrulewidth}{0.4pt}
\definecolor{codegray}{rgb}{0.45,0.45,0.45}
\lstdefinestyle{cppstyle}{basicstyle=\ttfamily\scriptsize,keywordstyle=\bfseries,commentstyle=\itshape\color{codegray},numbers=left,numberstyle=\tiny\color{codegray},numbersep=8pt,frame=single,framerule=0.4pt,rulecolor=\color{codegray},backgroundcolor=\color{gray!5},breaklines=true,tabsize=4,showstringspaces=false,language=C++}
\lstset{style=cppstyle}
\theoremstyle{plain}
\newtheorem{theorem}{Theorem}[chapter]
\newtheorem{lemma}[theorem]{Lemma}
\newtheorem{corollary}[theorem]{Corollary}
\newtheorem{proposition}[theorem]{Proposition}
\theoremstyle{definition}
\newtheorem{definition}[theorem]{Definition}
\newtheorem{example}[theorem]{Example}
\newtheorem{remark}[theorem]{Remark}
\theoremstyle{remark}
\newtheorem{exercise}[theorem]{Exercise}
\newtheorem{problem}[theorem]{Problem}
\newenvironment{cpplisting}{\begin{center}\small}{\end{center}}

\begin{document}

\chapter{Online Algorithms}
\label{ch:online-algorithms}

An \emph{online algorithm} must process its input piece by piece, making
irrevocable decisions at each step without knowledge of future requests.
This stands in sharp contrast to an \emph{offline algorithm}, which
receives the entire input in advance and may compute an optimal solution
by examining all of it.  The central question in the study of online
algorithms is: \emph{how much worse must an online algorithm perform
compared to an optimal offline algorithm that sees the entire request
sequence in advance?}

This degradation is captured by the \emph{competitive ratio}, defined as
the worst-case ratio of the cost incurred by the online algorithm to that
of the optimal offline algorithm.  Randomization plays a pivotal role in
online algorithms because it makes the algorithm's behavior unpredictable
to an adversary.  Whereas a deterministic online algorithm is completely
determined by the request sequence and hence can be ``foiled'' by a
carefully chosen worst-case input, a randomized algorithm can distribute
its cost over many possible actions, rendering the adversary's task of
exploiting worst-case structure significantly harder.

In this chapter we study several fundamental online problems, with
particular emphasis on the \emph{online paging problem} and the
\emph{$k$-server problem}.  We develop randomized algorithms that achieve
competitive ratios strictly better than what deterministic algorithms can
guarantee, and we analyze the interplay between different adversary models
that govern the power of an adversary against randomized algorithms.

%% ----------------------------------------------------------------
\section{The Online Paging Problem}
\label{sec:paging}
%% ----------------------------------------------------------------

\begin{definition}[Online paging]
Let $\mathcal{U}$ be a universe of $n$ pages and let the algorithm
maintain a \emph{cache} (or \emph{fast memory}) of $k$ pages, where
$k < n$.  A \emph{request sequence} $\sigma = \sigma_1, \sigma_2,
\ldots, \sigma_m$ is a sequence of pages drawn from $\mathcal{U}$.  When
request $\sigma_i$ arrives, the algorithm must ensure that $\sigma_i$ is
in the cache; if it is not, a \emph{cache miss} (or \emph{page fault})
occurs and the algorithm must bring $\sigma_i$ into the cache by evicting
some page currently stored there.  The cost of a request is 1 if it
causes a miss and 0 if it is a hit.  The total cost is the number of
misses incurred over the entire sequence.
\end{definition}

The goal is to minimize the total number of cache misses.  The
\emph{optimal offline algorithm} (denoted OPT) knows the entire request
sequence $\sigma$ in advance and computes the minimum possible number of
misses.  We say an algorithm $A$ is \emph{$c$-competitive} if for every
request sequence $\sigma$,
\[
  \mathrm{cost}_A(\sigma) \;\le\; c \cdot \mathrm{cost}_{\mathrm{OPT}}(\sigma) \;+\; O(1).
\]

\subsection{Deterministic Algorithms and Their Limitations}

Several natural deterministic strategies are immediately
available:

\begin{itemize}[nosep]
  \item \textbf{LRU (Least Recently Used):} Evict the page whose last
        use was furthest in the past.
  \item \textbf{FIFO (First-In, First-Out):} Evict the page that has
        been in the cache the longest.
  \item \textbf{FAR (Furthest in the Future):} Evict the page whose next
        use is furthest in the future (this is the optimal offline
        algorithm).
\end{itemize}

While FAR achieves the optimal cost, it requires knowledge of future
requests and is therefore offline.  LRU and FIFO are online, but both
can be forced to perform poorly.

\begin{example}\label{ex:paging-adversary}
Consider a cache of size $k=3$ with page universe $\{1, 2, 3, 4\}$.
Suppose the cache initially contains $\{1, 2, 3\}$.  Consider the
request sequence
\[
  \sigma = 4,\; 1,\; 4,\; 2,\; 4,\; 3,\; 4,\; 1,\; 4,\; 2,\; 4,\; 3,\; 4,\; \ldots
\]
In the subsequence $4, 1, 4, 2, 4, 3, 4$, every third request is page 4,
and pages 1, 2, 3 cycle through the remaining slots.  Every request
causes a miss for LRU (and similarly for FIFO), for a total cost of
approximately $\tfrac{4}{3} \cdot |\sigma|$.  However, the optimal
offline algorithm, seeing the pattern in advance, can keep page 4 in the
cache at all times and suffer only one miss per three requests, for a
total cost of approximately $\tfrac{1}{3} \cdot |\sigma|$.
\end{example}

\subsection{Lower Bound for Deterministic Algorithms}

The example above generalizes.  We now prove the classical lower bound
showing that no deterministic online paging algorithm can achieve a
competitive ratio better than $k-1$.

\begin{theorem}\label{thm:paging-det-lb}
Any deterministic online algorithm $A$ for $k$-page caching satisfies
\[
  \sup_\sigma \frac{\mathrm{cost}_A(\sigma)}{\mathrm{cost}_{\mathrm{OPT}}(\sigma)} \;\ge\; k - 1.
\]
\end{theorem}

\begin{proof}
Let the universe consist of $k+1$ pages $\{p_0, p_1, \ldots, p_k\}$.
Consider the following adversarial strategy.  The adversary presents
requests in blocks.  In each block, the adversary requests the $k+1$
pages in some order that avoids the page currently in position $k$ of
algorithm $A$'s internal ordering (i.e., the page $A$ would evict
next).  Specifically, the adversary maintains a set $S$ of $k+1$ pages.
At the start of each block, $A$'s cache contains some $k$-element subset
of $S$.  The adversary requests the one page of $S$ not in $A$'s cache,
causing a miss.  Then it continues requesting other pages of $S$ to
rearrange $A$'s cache without reducing its own cost.

More precisely, suppose at the start of a block, $A$'s cache contains
pages $q_1, q_2, \ldots, q_k \in S$ and the missing page is $p$.  The
adversary presents the single request $p$.  Algorithm $A$ must evict
some page $q_j$ to bring in $p$, incurring a miss.  Meanwhile, OPT can
serve this request with a hit if $p$ was already in OPT's cache.  By
choosing $p$ to be the unique page of $S$ not in $A$'s cache, the
adversary guarantees $A$ misses while OPT can serve the request with cost
0 (after the first block, OPT keeps all of $S \setminus \{q\}$ for an
appropriate page $q$ in its cache, using $k$ slots for $k$ pages).

After each block, $A$'s cache has changed (one page was evicted and $p$
was inserted), but still consists of $k$ pages from $S$.  The adversary
repeats, choosing the new missing page.  Each block costs $A$ at least 1
miss while OPT pays 0 (after the initial cost of loading its cache).  If
we group the sequence into rounds of $k$ blocks each, we can force
$A$ to miss $k-1$ times per round while OPT misses only once.  Hence
\[
  \frac{\mathrm{cost}_A(\sigma)}{\mathrm{cost}_{\mathrm{OPT}}(\sigma)} \;\ge\; k-1.
\]
\end{proof}

The above bound is tight for randomized algorithms, as we shall see.

%% ----------------------------------------------------------------
\section{Adversary Models}
\label{sec:adversary-models}
%% ----------------------------------------------------------------

The competitive analysis of randomized algorithms requires a careful
treatment of the adversary.  Unlike the deterministic case, where the
input is fixed and the algorithm is deterministic, the randomized
algorithm involves internal randomness, and the adversary's knowledge
of this randomness fundamentally affects the analysis.

\begin{definition}[Competitive ratio]
An online algorithm $A$ (possibly randomized) is
\emph{$C$-competitive} if, for every input sequence $\sigma$,
\[
  \mathbf{E}\!\bigl[\mathrm{cost}_A(\sigma)\bigr] \;\le\; C \cdot \mathrm{cost}_{\mathrm{OPT}}(\sigma) \;+\; O(1),
\]
where the expectation is over the random coins of $A$.
\end{definition}

\subsection{The Three Adversary Models}

There are three natural adversary models for randomized algorithms,
ordered from weakest to strongest:

\begin{enumerate}[nosep]
  \item \textbf{Oblivious adversary:} The adversary must fix the entire
        input sequence $\sigma$ \emph{before} the algorithm begins
        execution.  The adversary has no knowledge of the algorithm's
        random coins.  This is the \emph{weakest} adversary.

  \item \textbf{Adaptive online adversary:} The adversary produces
        requests one at a time.  When producing request $\sigma_i$, the
        adversary knows the algorithm's random coins and the algorithm's
        behavior up to step $i-1$, but does \emph{not} know the
        algorithm's future random coins (those that will be used after
        step $i$).  This adversary is \emph{intermediate} in power.

  \item \textbf{Adaptive offline adversary:} The adversary sees the
        algorithm's entire sequence of random coins (and hence the
        algorithm's complete behavior) and then produces the worst-case
        input sequence.  Equivalently, the adversary can view the
        algorithm as a function from random strings to behaviors, and
        choose the input to maximize cost against the worst-case random
        string.  This is the \emph{strongest} adversary.
\end{enumerate}

\begin{proposition}\label{prop:adversary-hierarchy}
For any randomized online algorithm $A$ and problem $\Pi$, if $A$ is
$C$-competitive against an adaptive offline adversary, then $A$ is also
$C$-competitive against an adaptive online adversary, and hence against an
oblivious adversary:
\[
  C_{\text{oblivious}} \;\le\; C_{\text{adaptive online}} \;\le\; C_{\text{adaptive offline}}.
\]
\end{proposition}

\begin{proof}
If $A$ is $C$-competitive against an adaptive offline adversary, then for
every input $\sigma$ and every random string $r$ of $A$,
$\mathrm{cost}_A(\sigma, r) \le C \cdot \mathrm{cost}_{\mathrm{OPT}}(\sigma) + O(1)$.
Taking expectations over $r$ gives the bound against oblivious
adversaries.  The adaptive online adversary is weaker than the adaptive
offline one because it must commit to each request without seeing the
algorithm's future random coins, which can only reduce the adversary's
ability to choose the worst-case input.
\end{proof}

The hierarchy is strict in general: there exist problems and algorithms
where the competitive ratio against an oblivious adversary is strictly
better than against an adaptive offline adversary.  However, as we shall
see in Section~\ref{sec:relating-adversaries}, for the paging problem the
situation is more nuanced.

\subsection{The Randomized Competitive Guarantee}

In the randomized setting, the competitive ratio is defined as the
worst-case (over input sequences) of the ratio of the expected cost of
the algorithm to the cost of OPT.  Formally:

\begin{definition}[Randomized competitive ratio]\label{def:rand-competitive}
A randomized algorithm $A$ is \emph{$C$-competitive against adversary
$\mathcal{A}$} if for every input sequence $\sigma$ (or every
$\sigma$ chosen by $\mathcal{A}$),
\[
  \mathbf{E}\!\bigl[\mathrm{cost}_A(\sigma)\bigr] \;\le\; C \cdot \mathrm{cost}_{\mathrm{OPT}}(\sigma) \;+\; O(1),
\]
where $\mathcal{O}(1)$ is a constant independent of $|\sigma|$.
\end{definition}

%% ----------------------------------------------------------------
\section{Paging against an Oblivious Adversary}
\label{sec:paging-oblivious}
%% ----------------------------------------------------------------

We now present a randomized algorithm for the $k$-page paging problem
that achieves a competitive ratio of $O(k)$ against an oblivious
adversary.

\subsection{The Randomized Marking Algorithm}

\begin{definition}[Marking algorithm]\label{def:marking}
The \emph{randomized marking algorithm} for $k$-page caching operates
as follows.  Each page in the cache carries a \emph{mark} (either
\emph{marked} or \emph{unmarked}).

\begin{itemize}[nosep]
  \item \textbf{On a request to page $p$:}
    \begin{enumerate}
      \item If $p$ is in the cache, mark $p$ (if not already marked).
            This is a \emph{hit}.
      \item If $p$ is not in the cache (a \emph{miss}):
        \begin{enumerate}
          \item If all $k$ pages in the cache are marked, unmark all
                pages and then mark $p$.  (Equivalently, start a new
                ``phase.'')
          \item Otherwise (at least one page in the cache is unmarked),
                \emph{evict a uniformly random unmarked page} and
                insert $p$, marking it.
        \end{enumerate}
    \end{enumerate}
\end{itemize}
\end{definition}

The key insight is that the marking algorithm partitions the request
sequence into \emph{phases}.  A new phase begins whenever all $k$ pages
in the cache are marked and a miss occurs.  Within a phase, the
algorithm only evicts randomly among unmarked pages.

\subsection{Analysis against an Oblivious Adversary}

We analyze the marking algorithm using a potential function argument.
Consider the configuration of the algorithm's cache $A$ and the optimal
offline algorithm's cache $O$ at any point during the execution.

\begin{definition}[Potential function]\label{def:paging-potential}
At any point in the execution, define the potential
\[
  \Phi \;=\; \bigl|\{p \in A \setminus O\}\bigr|,
\]
i.e., the number of pages in $A$'s cache that are not in $O$'s cache.
Here $A$ denotes the set of pages in the algorithm's cache and $O$
denotes the set of pages in OPT's cache.
\end{definition}

Note that $0 \le \Phi \le k$ since $|A| = |O| = k$.

\begin{lemma}\label{lem:marking-hit}
Within a phase, every request to a page that is in the cache is served
with cost 0 (a hit) by both the marking algorithm and OPT.
\end{lemma}

\begin{proof}
Within a phase, once a page is marked, it is never evicted (the algorithm
only evicts unmarked pages).  Thus all pages that have been requested at
least once during the current phase remain in the cache until the phase
ends.  OPT may choose to evict such a page, but it need not.
\end{proof}

\begin{lemma}\label{lem:marking-miss-prob}
When the marking algorithm suffers a miss on request $\sigma_i$ (with at
least one unmarked page in the cache), the probability that the
evicted page is one that is not in OPT's cache is at least $1/k$.
\end{lemma}

\begin{proof}
At the moment of eviction, the marking algorithm has $k$ pages in cache,
of which some subset $U$ of size $|U| \ge 1$ are unmarked.  It evicts a
page chosen uniformly at random from $U$.  Let $B = A \setminus O$ be the
set of pages in $A$'s cache but not in $O$'s cache.  Since $|A| = |O| =
k$, we have $|B| = |A \setminus O| = |O \setminus A|$.  At least $|B|$
pages in $A$ are ``bad'' (not in OPT).

The unmarked pages $U$ are a subset of $A$.  At most $k - |U|$ pages of
$A$ are marked.  In the worst case, all pages of $O \cap A$ could be
marked, so $|U \cap B| \ge |U| - |O \cap A \cap U^c|$ is not directly
computable.  However, observe that $|A \setminus O| = |B|$ and $|U|
\ge 1$.  Since $|A \cap O| \le k - |B|$ and $|U \cap A| = |U|$, we
have
\[
  |U \cap B| \;=\; |U| - |U \cap O| \;\ge\; |U| - |A \cap O| \;\ge\; |U| - (k - |B|).
\]
In the simplest case where $|U| = k$ (all pages are unmarked at the
start of a phase), we get $|U \cap B| = |B| \ge 1$ (if $|B| > 0$).
The probability of evicting a page in $B$ is $|U \cap B|/|U|$.

We use a cleaner argument: when a miss occurs, the newly arriving page
is not in $A$.  Either it is in $O$ or it is not.  Either way, $|A
\setminus O|$ does not decrease by more than 1 due to the insertion.
For the eviction step, note that the algorithm picks uniformly at random
from $U$.  Among $U$, at least one page must be in $A \setminus O$
whenever $|A \setminus O| > 0$, because $|U| \le k$ and $|A \cap O| \le
k - 1$ (since the new page is not in $A$, so $|A \cap O|$ could be at
most $k$).  In fact, since $|A \setminus O|$ counts pages in $A$ not in
$O$, and $|U| \ge 1$, by a counting argument at least a $1/k$ fraction
of $U$ lies in $A \setminus O$:
\[
  \frac{|U \cap (A \setminus O)|}{|U|} \;\ge\; \frac{|A \setminus O| - (k - |U|)}{|U|}.
\]
When $|A \setminus O| \ge 1$, this is at least $1/k$ because in the worst
case $|A \setminus O| = 1$ and $|U| = k$, giving probability $1/k$.
\end{proof}

We now state the main theorem.

\begin{theorem}\label{thm:marking-2k}
The randomized marking algorithm is $2k$-competitive against an oblivious
adversary for the $k$-page paging problem.
\end{theorem}

\begin{proof}
We use the potential function $\Phi = |A \setminus O|$ as defined in
Definition~\ref{def:paging-potential}.  Recall that $0 \le \Phi \le k$.

We partition the request sequence into \emph{phases}.  A phase ends when
all $k$ pages in the cache are marked and a new miss occurs (triggering
a global unmark).  Consider one phase.  Let $m_A$ be the number of
misses the marking algorithm suffers in this phase, and let $m_O$ be the
number of misses OPT suffers in this phase (for the requests in this
phase only, considering OPT's cache state at the start of the phase).

\textbf{Claim:} $m_O \ge m_A - k$.

\emph{Proof of claim:} Each miss by the marking algorithm brings a new
page into the cache.  At the end of the phase, all $k$ pages in the
cache are marked, meaning $k$ distinct pages were brought in during the
phase (each was requested at least once while unmarked).  OPT starts
the phase with $k$ pages.  Each miss by OPT introduces one new page.
To serve all the requests in the phase that cause misses for the marking
algorithm, OPT must also miss on at least $m_A - k$ of them, because
OPT's cache can contain at most $k$ of the $m_A$ pages that the marking
algorithm brought in (OPT might already have some of them, but not all
$m_A$).  Hence $m_O \ge m_A - k$.

Now we do the potential function analysis.  Over a phase, define
$\Delta\Phi$ as the change in $\Phi$ from the start to the end of the
phase.  Since $\Phi$ measures $|A \setminus O|$ and the phase ends with
a new marking configuration, we have $\Delta\Phi \le k$ (since $\Phi \le
k$ at all times).

We relate costs and potential.  Consider the amortized cost of the
marking algorithm per phase:
\[
  m_A + \mathbf{E}[\Delta\Phi] \;\le\; m_A + k.
\]
By the claim, $m_A \le m_O + k$.  Thus the expected amortized cost is at
most $m_O + 2k$.  Since $\Phi \ge 0$ initially and $\Phi \ge 0$ always,
summing over all phases gives
\[
  \mathbf{E}[\mathrm{cost}_A(\sigma)] \;=\; \sum_{\text{phases}} m_A \;\le\; \sum_{\text{phases}} (m_O + 2k) \;=\; \mathrm{cost}_{\mathrm{OPT}}(\sigma) + 2k \cdot (\text{number of phases}).
\]
Each phase has at least one miss by the marking algorithm (otherwise the
phase wouldn't end), and OPT misses at least once per phase (since the
page requested at the phase boundary is not in $A$'s cache and hence
the adversary can arrange for OPT to miss as well, or more precisely the
number of phases is at most $\mathrm{cost}_{\mathrm{OPT}}(\sigma)$ since
each phase boundary corresponds to an OPT miss).  Hence the number of
phases is at most $\mathrm{cost}_{\mathrm{OPT}}(\sigma)$, and
\[
  \mathbf{E}[\mathrm{cost}_A(\sigma)] \;\le\; \mathrm{cost}_{\mathrm{OPT}}(\sigma) + 2k \cdot \mathrm{cost}_{\mathrm{OPT}}(\sigma) \;=\; (2k+1) \cdot \mathrm{cost}_{\mathrm{OPT}}(\sigma).
\]
This gives a competitive ratio of $2k+1$, which is $O(k)$.
\end{proof}

\subsection{Improvement to $k$-Competitive}

The bound above can be tightened to a $k$-competitive algorithm using a
more refined analysis.  The key observation is that the potential
function $\Phi$ can be defined more carefully.

\begin{theorem}\label{thm:marking-k}
There exists a randomized marking algorithm that is $k$-competitive
against an oblivious adversary.
\end{theorem}

\begin{proof}[Proof sketch]
Consider the same marking algorithm but with a refined potential function.
Define $\Phi = |A \setminus O|$ as before.  Analyze the algorithm on a
per-request basis rather than per-phase.  For a request that is a hit
for the marking algorithm, the cost is 0 and $\Phi$ does not change (the
page was already in $A$ and marked).  For a request that is a miss:

\begin{itemize}[nosep]
  \item The marking algorithm pays cost 1 and brings in a new page
        $\sigma_i$.  This changes $|A \setminus O|$ by at most $+1$ (if
        $\sigma_i \notin O$) or $-1$ (if $\sigma_i \in O$).
  \item The eviction of a random unmarked page $q$ changes $|A \setminus
        O|$ by $-1$ if $q \in A \setminus O$ and $0$ if $q \in A \cap
        O$.
\end{itemize}

For a miss where $|A \setminus O| > 0$, the expected change in potential
due to the eviction is at most $-1/k$ (by the argument that at least a
$1/k$ fraction of unmarked pages are in $A \setminus O$).  Combining
the cost of the miss with the expected change in potential, we get an
amortized cost of at most $1 - 1/k$ per miss.  Summing over all
misses and using $0 \le \Phi \le k$ gives
\[
  \mathbf{E}[\mathrm{cost}_A(\sigma)] + 0 \;\le\; (1 - 1/k)^{-1} \cdot \mathrm{cost}_{\mathrm{OPT}}(\sigma) + k \;=\; \frac{k}{k-1} \cdot \mathrm{cost}_{\mathrm{OPT}}(\sigma) + k,
\]
which is $k/(k-1)$-competitive.  However, a more careful analysis that
considers the interaction between the phases and the potential yields a
$k$-competitive bound.  The full details involve tracking the potential
not just as $|A \setminus O|$ but weighted by a function of the number
of marked pages.
\end{proof}

%% ----------------------------------------------------------------
\section{Relating the Adversary Models}
\label{sec:relating-adversaries}
%% ----------------------------------------------------------------

A natural question is whether the choice of adversary model affects the
competitive ratio achievable by randomized algorithms for paging.  We
have seen that $k$-competitive algorithms exist against oblivious
adversaries.  What about against stronger adversaries?

\begin{theorem}\label{thm:paging-adversary-equivalence}
For the $k$-page paging problem, any randomized algorithm that is
$C$-competitive against an adaptive offline adversary is also
$C$-competitive against an adaptive online adversary.
\end{theorem}

\begin{proof}[Proof sketch]
Suppose $A$ is $C$-competitive against an adaptive offline adversary.
We show $A$ is also $C$-competitive against an adaptive online adversary
$\mathcal{A}_{\mathrm{on}}$.

The adaptive online adversary $\mathcal{A}_{\mathrm{on}}$ produces
requests one at a time.  At step $i$, $\mathcal{A}_{\mathrm{on}}$
chooses $\sigma_i$ based on the algorithm's coins $r$ and the algorithm's
behavior up to step $i-1$.  Consider the adaptive offline adversary
$\mathcal{A}_{\mathrm{off}}$ that, given the algorithm's entire random
string $r$, simulates $\mathcal{A}_{\mathrm{on}}$ by running the same
adaptive strategy.  Since $\mathcal{A}_{\mathrm{off}}$ knows $r$ in
advance, it can predict exactly what $\mathcal{A}_{\mathrm{on}}$ would
do at each step.  Thus for every $r$, the sequence produced by
$\mathcal{A}_{\mathrm{off}}$ is identical to that produced by
$\mathcal{A}_{\mathrm{on}}$ against $A$ with coins $r$.  Therefore
\[
  \mathbf{E}_r\!\bigl[\mathrm{cost}_A(r, \sigma_r)\bigr] \;\le\; C \cdot \mathrm{cost}_{\mathrm{OPT}}(\sigma_r) + O(1)
\]
for every $r$, where $\sigma_r$ is the sequence produced by
$\mathcal{A}_{\mathrm{on}}$ when $A$ uses coins $r$.  Taking
expectations over $r$ yields the $C$-competitive bound against
$\mathcal{A}_{\mathrm{on}}$.
\end{proof}

\begin{remark}
The above argument shows that the adaptive online adversary is no
stronger than the adaptive offline adversary for \emph{any} online
problem, not just paging.  The converse is not true in general: there
exist problems where the adaptive offline adversary is strictly stronger.
The hierarchy collapses for paging because the adaptive online adversary
can be ``simulated'' by the offline adversary.
\end{remark}

\begin{theorem}\label{thm:paging-rand-lb}
No randomized algorithm for $k$-page caching can achieve a competitive
ratio better than $k$ against an adaptive offline adversary (and hence
against any adversary).
\end{theorem}

\begin{proof}
The argument is an adversary strategy.  Consider a universe of $k+1$
pages.  The adversary maintains a set $S$ of $k+1$ candidate pages.
At each step, the adversary requests a page that is not in the
algorithm's cache (if possible).  Since the algorithm's cache has $k$
pages and $|S| = k+1$, there is always at least one page in $S$ not in
the algorithm's cache.

The adversary can use the algorithm's random coins to predict which page
the algorithm will evict and choose the request to force a miss.
Specifically, given the algorithm's coins $r$, the adversary knows the
algorithm's complete strategy.  It requests the page of $S$ not in the
algorithm's cache.  The algorithm must miss (cost 1).  Meanwhile, OPT
can always serve this request with cost at most $1/k$ on average
(amortized), because OPT can rotate its cache contents to keep $k$ of
the $k+1$ pages, missing only once every $k$ requests.  Thus each
request costs the algorithm 1 and OPT at most $1/k$, giving a ratio of
at least $k$.

More precisely, the adversary requests pages from $S$ in a round-robin
fashion, always picking the one not in the algorithm's cache.  Each
round of $k$ requests costs the algorithm $k$ misses while OPT misses
at most once (after an initial cost of 1 to load its cache).  Hence
\[
  \frac{\mathrm{cost}_A(\sigma)}{\mathrm{cost}_{\mathrm{OPT}}(\sigma)} \;\ge\; k.
\]
This holds for every random string $r$, so the expectation is at least
$k$ as well.
\end{proof}

Combining with Theorem~\ref{thm:marking-k}, we see that the
$k$-competitive ratio is \emph{optimal} for randomized paging against
any adversary.

%% ----------------------------------------------------------------
\section{The Adaptive Online Adversary}
\label{sec:adaptive-online}
%% ----------------------------------------------------------------

The adaptive online adversary model provides a natural middle ground
between the oblivious and fully adaptive offline models.  In this
section, we establish lower bounds and discuss the implications of this
model for paging and related problems.

\begin{theorem}\label{thm:adaptive-online-lb}
No randomized algorithm for $k$-page caching can achieve a competitive
ratio better than $k$ against an adaptive online adversary.
\end{theorem}

\begin{proof}
The proof is essentially the same as that of
Theorem~\ref{thm:paging-rand-lb}, adapted to the online setting.  The
adaptive online adversary produces requests one at a time.  At each
step, the adversary knows the algorithm's coins $r$ used so far (those
that affected the algorithm's behavior up to this point) and the
algorithm's current cache state.

The adversary maintains a set $S$ of $k+1$ pages.  At each step, the
adversary selects $\sigma_i$ to be the page of $S$ not in the
algorithm's current cache.  Since $|S| = k+1$ and the cache holds $k$
pages, such a page always exists.  This forces a miss (cost 1).

Meanwhile, OPT can keep $k$ pages of $S$ in its cache and rotate them.
In the worst case, OPT misses once every $k$ requests.  Hence the
cost ratio per request is at least $k$.

The adaptive online adversary does not need to know the algorithm's
future random coins to execute this strategy---it only needs to know the
current cache state, which is determined by the algorithm's past coins
and requests.  Therefore this strategy is valid against the adaptive
online adversary.
\end{proof}

\subsection{Comparison of Models for Paging}

For the paging problem, the competitive ratios achievable against the
three adversary models are:

\begin{center}
\begin{tabular}{lcc}
\toprule
\textbf{Adversary model} & \textbf{Best deterministic} & \textbf{Best randomized} \\
\midrule
Oblivious          & $k-1$ & $k$ \\
Adaptive online    & $k-1$ & $k$ \\
Adaptive offline   & $k-1$ & $k$ \\
\bottomrule
\end{tabular}
\end{center}

The gap between the best deterministic ratio ($k-1$) and the best
randomized ratio ($k$) is small for paging.  The advantage of
randomization for paging is primarily in the \emph{constant} hidden in
the $O(\cdot)$ notation and in the adaptability of the algorithm.

\begin{remark}
For other problems, such as the $k$-server problem discussed in the next
section, the gap between adversary models is much more dramatic, and
randomization provides a significantly larger advantage.
\end{remark}

%% ----------------------------------------------------------------
\section{The $k$-Server Problem}
\label{sec:k-server}
%% ----------------------------------------------------------------

The $k$-server problem is a fundamental generalization of paging and one
of the most important problems in the theory of online algorithms.

\begin{definition}[$k$-server problem]\label{def:k-server}
Let $(\mathcal{M}, d)$ be a metric space with $n$ points.  There are $k$
servers located at points of $\mathcal{M}$ (initially at some fixed
positions $s_1, \ldots, s_k$).  Requests arrive one at a time: request
$r_t$ is a point in $\mathcal{M}$.  When $r_t$ arrives, some server must
be moved to $r_t$.  If a server is already at $r_t$, the cost is 0.
Otherwise, if server $s_j$ at position $p_j$ is moved to $r_t$, the
cost is $d(p_j, r_t)$.  The goal is to minimize the total cost of all
movements.
\end{definition}

Note that when $k=1$ and $\mathcal{M}$ is a set of pages with the
discrete metric, the $k$-server problem reduces to a variant of the
paging problem.  In general, the $k$-server problem captures many
important applications, including caching, load balancing, and
data migration.

\subsection{Deterministic Lower Bound}

\begin{theorem}[Manasse, McGeoch, and Sleator]\label{thm:k-server-det-lb}
Any deterministic online algorithm for the $k$-server problem on an
$n$-point metric space is at least $k$-competitive.  More precisely,
there exist metric spaces on which no deterministic algorithm achieves a
competitive ratio better than $k$.
\end{theorem}

\begin{proof}[Proof sketch]
Consider the metric space consisting of $k+1$ points arranged in a
line (or a complete graph with equal edge weights).  The adversary
places $k$ servers at the first $k$ points and then alternates
requesting the $(k+1)$-th point and one of the first $k$ points.  At
each step, the adversary can force the algorithm to move a server across
a long distance, while OPT can move a server only a short distance.
The resulting competitive ratio is at least $k$.
\end{proof}

\subsection{Randomized $k$-Competitive Algorithm}

\begin{theorem}[Koutsoupias and Papadimitriou]\label{thm:k-server-kp}
There exists a randomized online algorithm for the $k$-server problem
that is $k$-competitive on every metric space.
\end{theorem}

The Koutsoupias--Papadimitriou (KP) algorithm is a landmark result.  We
describe the algorithm and its analysis.

\subsubsection{The Algorithm}

The KP algorithm maintains, in addition to the actual server
configuration, a probability distribution over all possible
configurations.  At each step, it moves servers to minimize a potential
function $\Phi$ defined over pairs of configurations.

\begin{definition}[KP potential function]\label{def:kp-potential}
Let $S = (s_1, \ldots, s_k)$ be the algorithm's server configuration
and let $\mathcal{C}$ be the set of all possible configurations of $k$
servers in $\mathcal{M}$.  For configurations $C, C' \in \mathcal{C}$,
define
\[
  d(C, C') \;=\; \min_{\pi} \sum_{i=1}^{k} d\bigl(s_i, s'_{\pi(i)}\bigr),
\]
the minimum-cost matching between the servers of $C$ and $C'$.  The
potential function is
\[
  \Phi(S, R) \;=\; \max_{C \in \mathcal{C}} \bigl[ d(S, C) - \alpha \cdot d(R, C) \bigr],
\]
where $R$ is the configuration of OPT and $\alpha = k$ is the
competitive ratio.  Here $d(S, C)$ is the minimum-cost matching
between the algorithm's servers $S$ and the servers in configuration $C$.
\end{definition}

The algorithm proceeds as follows: upon a request at point $r$, the
algorithm moves a server to $r$ so as to minimize $\Phi$.  It can be
shown that a greedy minimization of $\Phi$ (moving the server that
minimizes the maximum over $C$ of the change in $d(S,C) - k \cdot
d(R,C)$) achieves a competitive ratio of $k$.

\subsubsection{Analysis}

\begin{lemma}\label{lem:kp-step}
At each step, the expected change in the potential $\Phi$ satisfies
\[
  \mathbf{E}\!\bigl[\Delta\Phi\bigr] \;\le\; \alpha \cdot \mathrm{cost}_{\mathrm{OPT}}(r_t) - \mathrm{cost}_A(r_t),
\]
where $\alpha = k$.
\end{lemma}

\begin{proof}[Proof sketch]
When request $r_t$ arrives, OPT may move a server to $r_t$ at cost
$\mathrm{cost}_{\mathrm{OPT}}(r_t) \le d(s_j, r_t)$ for some server
$s_j$ of OPT.  The algorithm chooses a server $s_i$ to move to $r_t$ so
that the decrease in $\mathrm{cost}_A(r_t)$ is at least $1/\alpha$ times
the decrease in $\Phi$.  By the definition of $\Phi$ as a maximum over
configurations $C$, the change $\Delta\Phi$ can be bounded using the
triangle inequality on the metric space.  The key inequality is
\[
  \mathrm{cost}_A(r_t) + \mathbf{E}[\Delta\Phi] \;\le\; \alpha \cdot \mathrm{cost}_{\mathrm{OPT}}(r_t),
\]
which, when summed over all steps and using $0 \le \Phi \le k \cdot
\mathrm{diam}(\mathcal{M})$ (which is bounded), gives the $k$-competitive
ratio.
\end{proof}

\begin{proof}[Proof of Theorem~\ref{thm:k-server-kp}]
By Lemma~\ref{lem:kp-step}, at each step $t$,
\[
  \mathrm{cost}_A(r_t) + \mathbf{E}[\Delta\Phi_t] \;\le\; k \cdot \mathrm{cost}_{\mathrm{OPT}}(r_t).
\]
Summing over all $T$ steps:
\[
  \sum_{t=1}^{T} \mathrm{cost}_A(r_t) + \mathbf{E}[\Phi_T - \Phi_0] \;\le\; k \sum_{t=1}^{T} \mathrm{cost}_{\mathrm{OPT}}(r_t).
\]
Since $\Phi_0$ is determined by the initial configurations and $\Phi_T \ge
0$, and both are bounded by $k$ times the diameter of the metric space,
we obtain
\[
  \mathrm{cost}_A(\sigma) \;\le\; k \cdot \mathrm{cost}_{\mathrm{OPT}}(\sigma) + \Phi_0 \;\le\; k \cdot \mathrm{cost}_{\mathrm{OPT}}(\sigma) + O(k \cdot \mathrm{diam}(\mathcal{M})).
\]
This gives the $k$-competitive ratio.
\end{proof}

\subsection{Connection to Paging}

\begin{corollary}\label{cor:k-server-paging}
When the metric space is the discrete metric on a set of $n$ pages and
$k$ is the cache size, the $k$-server problem generalizes the $k$-page
paging problem.  The KP algorithm achieves a $k$-competitive ratio for
paging, matching the lower bound of Theorem~\ref{thm:paging-rand-lb}.
\end{corollary}

\subsection{The Hirschberg--Sinclair Technique}

The analysis of the KP algorithm draws on ideas reminiscent of the
Hirschberg--Sinclair technique for the {\sc Load Balancing} problem.
The key idea is to consider \emph{all possible} optimal configurations
and use a potential function that is a maximum over these configurations.
This ``max-over-all-possibilities'' approach is a powerful technique in
the analysis of randomized online algorithms.

\begin{remark}
The $k$-competitive bound for $k$-server was the best known for many
years.  In 2021, Bubeck, Cohen, Lee, and Lee achieved a breakthrough
result showing that the randomized competitive ratio for $k$-server on
general metric spaces is at most $O(\log^2 k \log n)$, where $n$ is the
number of points in the metric space.  However, the $k$-competitive
bound remains the tight answer for many important special cases, and the
KP algorithm remains one of the most elegant constructions in the field.
\end{remark}

%% ----------------------------------------------------------------
%% C++ Implementation
%% ----------------------------------------------------------------

\subsection*{C++ Implementation}

\lstinputlisting[
  language=C++,
  caption={Online paging: LRU, FIFO, random, marking algorithm, and OPT.},
  label={lst:paging}]{src/chapter13/paging.h}

\lstinputlisting[
  language=C++,
  caption={$k$-server: greedy, randomized, and optimal offline algorithms.},
  label={lst:kserver}]{src/chapter13/k_server.h}

\lstinputlisting[
  language=C++,
  caption={Adversary simulation: oblivious and adaptive models.},
  label={lst:adversary}]{src/chapter13/adversary.h}

%% ----------------------------------------------------------------
\section*{Notes}
\label{sec:online-notes}
%% ----------------------------------------------------------------

The study of online algorithms was initiated by the work of
Sleator and Tarjan~\cite{sleator1985amortized}, who introduced the
competitive analysis framework and proved the $k-1$ lower bound for
deterministic paging (Theorem~\ref{thm:paging-det-lb}) and the
$O(k)$-competitive randomized marking algorithm
(Theorems~\ref{thm:marking-2k} and~\ref{thm:marking-k}).  The
marking algorithm was further analyzed by Karlin, Manasse, McGeoch,
and Sleator~\cite{karlin1988competitve}, who established the connection
between the marking algorithm and competitive ratios against oblivious
adversaries.

The Koutsoupias--Papadimitriou randomized $k$-competitive algorithm for
$k$-server~\cite{koutsoupias1995k} (Theorem~\ref{thm:k-server-kp}) was
a major open problem for nearly two decades.  The adversary hierarchy for
randomized online algorithms was studied extensively by Borodin and
El-Yaniv~\cite{borodin1998online}, who formalized the three adversary
models discussed in Section~\ref{sec:adversary-models} and proved the
hierarchy results of Proposition~\ref{prop:adversary-hierarchy} and
Theorem~\ref{thm:paging-adversary-equivalence}.

The lower bound for deterministic $k$-server
(Theorem~\ref{thm:k-server-det-lb}) is due to Manasse, McGeoch, and
Sleator~\cite{manasse1988greedy}.  The relationship between the
adversary models and the optimal competitive ratios was clarified by
Fiat and Mendel~\cite{fiat1997optimal}, who also obtained improved
algorithms for special metric spaces.

The recent breakthrough by Bubeck, Cohen, Lee, and
Lee~\cite{bubeck2021randomized} significantly improved the upper bound
for randomized $k$-server on general metric spaces, though the exact
competitive ratio remains an open problem.

Further reading can be found in the textbook by
Borodin and El-Yaniv~\cite{borodin1998online} and the survey by
Albers~\cite{albers2003new}.

%% ----------------------------------------------------------------
\section*{Problems}
\label{sec:online-problems}
%% ----------------------------------------------------------------

\begin{problem}\label{prob:lru-fifo}
Prove that the LRU and FIFO algorithms for $k$-page caching are
$\Theta(k)$-competitive against an oblivious adversary.  That is,
establish matching upper and lower bounds on their competitive ratios.
\end{problem}

\begin{problem}\label{prob:randomized-lb-oblivious}
Prove that no randomized algorithm for $k$-page caching can be better
than $k/(k-1)$-competitive against an oblivious adversary.
\end{problem}

\begin{problem}\label{prob:weighted-paging}
Consider the \emph{weighted paging} problem, where each page $p$ has a
weight $w_p > 0$ and the cost of a miss on page $p$ is $w_p$.  Design a
randomized algorithm that is $O(k \cdot W_{\max}/W_{\min})$-competitive
against an oblivious adversary, where $W_{\max}$ and $W_{\min}$ are the
maximum and minimum page weights.
\end{problem}

\begin{problem}\label{prob:two-server}
Analyze the $2$-server problem on a metric space of three points forming
an equilateral triangle with side length 1.  Show that the deterministic
lower bound is $2$ and describe a randomized $2$-competitive algorithm.
\end{problem}

\begin{problem}\label{prob:load-balancing}
The \emph{online load balancing} problem can be modeled as a special
case of the $k$-server problem.  Consider $k$ machines and jobs arriving
online.  Each job $j$ has a processing time $p_j$ on each machine, and
the cost of assigning job $j$ to machine $i$ is $p_{ij}$.  Show that
the $k$-server lower bound implies a lower bound of $k$ for competitive
deterministic load balancing.
\end{problem}

\begin{problem}\label{prob:oblivious-vs-adaptive}
Construct a problem (other than paging) where the competitive ratio
against an oblivious adversary is strictly better than against an
adaptive offline adversary.  Prove that the two ratios differ.
\end{problem}

\begin{problem}\label{prob:marking-analysis}
Complete the proof of Theorem~\ref{thm:marking-k} by providing the full
potential function argument showing that the randomized marking algorithm
is $k$-competitive against an oblivious adversary.  In particular,
explicitly define the refined potential function and verify that the
amortized cost per request is at most $k/(k-1) \cdot
\mathrm{cost}_{\mathrm{OPT}}$.
\end{problem}

\begin{problem}\label{prob:k-server-metric}
Prove that the $k$-server competitive ratio cannot depend on the number
of points $n$ in the metric space when $k$ is fixed.  That is, show
that for any fixed $k$, there exists a constant $C_k$ such that the
$k$-server problem on any $n$-point metric space admits a $C_k$-competitive
deterministic algorithm.
\end{problem}

\begin{problem}\label{prob:bit-packing}
Consider the \emph{bit} or \emph{pennies} online problem: at each step,
one of $n$ coins is placed on the table with heads up, and the
algorithm must flip one coin to tails (or declare a fault).  The goal is
to keep the number of heads-up coins below some threshold.  Formulate
this as an online algorithm problem and derive a competitive ratio using
techniques from this chapter.
\end{problem}

\begin{problem}\label{prob:kp-potential}
For the Koutsoupias--Papadimitriou algorithm
(Theorem~\ref{thm:k-server-kp}), prove Lemma~\ref{lem:kp-step} in full
detail.  Specifically, verify that the potential $\Phi$ is bounded, that
the greedy minimization step achieves the required inequality, and that
the final sum yields the $k$-competitive bound.
\end{problem}

\end{document}
